% ---------------------------------------------------------------
\section{Journal--institution citation networks}
\label{sec:networks}
% ---------------------------------------------------------------

%\paragraph{Corpus construction.}
We begin by selecting $N_s$ sources $\mathfrak{S}$ in a field, discipline or category.
Then we select $N_w$ works $\mathfrak{W}$ published in $\mathfrak{S}$ within the census
window years $t_c = [t^c_0 \ldots t^c_1]$ or target years $t_t = [t^t_0 \ldots t^t_1]$.
From the reference lists of the set of census works $\mathfrak{W}^c$ we identify the set of target works $\mathfrak{W}^t$ and form the filtered reference set $\mathfrak{R}=\{(i,j):w_i\in\mathfrak{W}^c,\,w_j\in\mathfrak{W}^t\}$.
Author affiliations in $\mathfrak{W}$ establish a set of authors $\mathfrak{A}'$ and institutions $\mathfrak{U}'$.
Let $\bar{w}_u$ be the mean annual work count for institution $u$.
We determine a lower limit $\tau_U$ on $\bar{w}_u$ by examining the work-institution frequency
distribution and define $\mathfrak{U}\subseteq\mathfrak{U}'$ as the institutions with
$\bar{w}_u\geq\tau_U$, with corresponding author set $\mathfrak{A}\subseteq\mathfrak{A}'$.
Works that lose all their authorships under this filter are dropped.
Each retained work $w_i$ has an \textit{institution weight} under author-fractional counting
\citep{waltmanFieldnormalizedCitationImpact2015},
\begin{equation}
  \omega_{i u} \;=\;
  \sum_{\ell:\,u\in\mathcal{u}_{i\ell}}
  \frac{1}{a_i}\,\frac{1}{u_{i\ell}},
  \label{eq:omega}
\end{equation}
where $\ell$ indexes the retained authors of work $w_i$, $a_i$ is their count, and $u_{i\ell}$ the number of
retained institutions of author $\ell$, so that $\sum_u\omega_{iu}=1$.

When a work $w_i$ references another work $w_j$ we say that $w_i$ gives attention to $w_j$. The attention given by a work can be either (a) the number of references, or (b) a fixed amount \citep{zittModifyingJournalImpact2008}. Let $R_i=|\{j:(i,j)\in\mathfrak{R}\}|$ be the reference count of work $w_i$ and $\bar{R}=|\mathfrak{R}|/N_w$ the mean reference count over $\mathfrak{W}$. We use the \textit{reference normalisation weight} $\rho_i$ to select case (a) or case (b) respectively:
\begin{equation}
	\rho_i \;=\;
	\begin{cases}
		 1 & \text{(full count)} \\
		\bar{R}/R_i & \text{(fixed count).}
	\end{cases}
	\label{eq:rho}
\end{equation}

We construct the set of \emph{units} [sources or institutions; \cite{pinskiCitationInfluenceJournal1976}]
$\mathfrak{P}=\mathfrak{S}\cup\mathfrak{U}$ with $N=N_s+N_u$ elements, ordered so that
sources occupy indices $1,\dots,N_s$ and institutions indices $N_s+1,\dots,N$.
A reference $(i,j) \in \mathfrak{R}$ gives attention across unit-type pairs:
source--source ($SS$), source--institution ($SI$), institution--source ($IS$), and
institution--institution ($II$).

For each work $w_i$ we define the source membership one-hot vector
$\mathbf{b}^S_i\in\mathbb{R}^{N_s}$ with $[\mathbf{b}^S_i]_s=\mathbf{1}[w_i\in s]$,
and the institution weight distribution vector
$\mathbf{b}^U_i\in\mathbb{R}^{N_u}$ with $[\mathbf{b}^U_i]_u=\omega_{iu}$
from \eqref{eq:omega}.
Attention flow between pairs of unit types is accumulated as four blocks over $\mathfrak{R}$ as rank-1 outer products,
\begin{equation}
	\mathbf{C}_{XY} \;=\; \sum_{(i,j)\,\in\,\mathfrak{R}}
	\rho_i\;\mathbf{b}^X_i\,\bigl(\mathbf{b}^Y_j\bigr)^\top,
	\qquad XY\in\{SS,\,SI,\,IS,\,II\},
	\label{eq:Cblock}
\end{equation}
with $\mathbf{b}^X_i$ weighting the referencing (citing) side and $\mathbf{b}^Y_j$ the referenced (cited) side.

The citation matrix assembles the four blocks:
\begin{equation}
	\mathbf{C}(\chi,\mathbf{m},\rho) \;=\;
	\begin{pmatrix}
		m_{SS}(1-\chi)^2\,\mathbf{C}_{SS} &
		m_{SI}\,\chi(1-\chi)\,\mathbf{C}_{SI} \\[4pt]
		m_{IS}\,\chi(1-\chi)\,\mathbf{C}_{IS} &
		m_{II}\,\chi^2\,\mathbf{C}_{II}
	\end{pmatrix},
	\label{eq:C}
\end{equation}
where the binary \textit{block mask} $\mathbf{m}=(m_{SS},m_{SI},m_{IS},m_{II})\in\{0,1\}^4$
is bibliometrically meaningful for
source-only $(1,0,0,0)$, institution-only $(0,0,0,1)$,
bipartite $(0,1,1,0)$, and full-joint $(1,1,1,1)$ combined blocks.
The mixing weight $\chi$ controls the source--institution balance
in the bipartite and full-joint cases and is inert in the single-mode cases.
The $\chi$ scaling preserves total citation flow across all block combinations,
since $(1-\chi)^2+2\chi(1-\chi)+\chi^2=1$.
We could name $\mathbf{C}$ the reference, attention or citation matrix and prefer the latter.

The model parameters are summarised in Table \ref{tab:params}. Note that the corpus excludes (a) census work references not giving attention to  target works, (b) non-census works giving attention to target works, and (c) authorships that include an omitted institution.

\begin{table}[h]
	\centering
	\caption{Model parameters $\boldsymbol{\theta}$.}
	\label{tab:params}
	\begin{tabularx}{\textwidth}{lllX}
		\toprule
		Symbol & Name & Type & Role \\
		\midrule
		\multicolumn{4}{@{}l@{}}{\textit{Corpus construction}} \\[2pt]
		$F\in\{E,B,A\}$          & field subset          & category        & restrict $\mathfrak{S}$ to Economics ($E$), Business ($B$), or all OAS ($A$); see Table~\ref{tab:windows} \\
		$t_x\in\{1,\ldots,7\}$  & time window           & integer         & census and target year ranges; see Table~\ref{tab:windows} \\
		$\tau_U$                 & institution threshold & integer         & minimum mean annual works to retain an institution \\[4pt]
		\multicolumn{4}{@{}l@{}}{\textit{Matrix construction}} \\[2pt]
		$\rho\in\{0,1\}$         & reference normalisation & binary        & 0: a work contributes $\bar{R}$ units \newline 1: a work contributes $R_i$ units \\
		$\mathbf{m}\in\{0,1\}^4$ & block mask            & binary 4-vector & select active unit-type pairs ($SS$, $SI$, $IS$, $II$) \\
		$\chi\in[0,1]$           & mixing weight         & continuous      & source--institution balance within each work \\[4pt]
		\multicolumn{4}{@{}l@{}}{\textit{Ranking}} \\[2pt]
		$\alpha\in[0,1)$         & damping factor        & continuous      & controls propagation through citation network \newline \hspace*{1em} ensures convergence \\
		\bottomrule
\end{tabularx}
\end{table}

\paragraph{Spectral ranking.}
Let $r_p=\sum_q C_{pq}$ be the total outgoing citation flow from unit $p$,
$\mathbf{D}_r=\mathrm{diag}(r_p)$ the corresponding diagonal matrix.
If $r_p=0$ (a dangling unit with no outgoing citations) the corresponding
row of $\mathbf{C}$ is zero and $\mathbf{H}$ is defined to have a zero row
there; the boundary condition below then supplies the necessary probability
mass.
Let $\boldsymbol{\mu}$ be an initial preference, or boundary condition
\citep{hubbellInputoutputApproachClique1965, vignaSpectralRanking2016},
with $\mu_p = 1/N$ for all $p\in\mathfrak{P}$.
The row-stochastic citation matrix is
$\mathbf{H}=\mathbf{D}_r^{-1}\mathbf{C}$.

The spectral ranking $\boldsymbol{\pi}$
\citep[see][for an overview]{vignaSpectralRanking2016}
is defined as the solution to the damped eigenvector problem
\begin{equation}
    \boldsymbol{\pi}
    \;=\;
    \alpha\mathbf{H}^\top\boldsymbol{\pi}
    +(1-\alpha)\boldsymbol{\mu},
    \qquad
    \mathbf{1}^\top\boldsymbol{\pi}=1,
    \quad
    \boldsymbol{\pi}\ge 0,
    \label{eq:spectral}
\end{equation}
with damping factor $\alpha\in[0,1)$.
The prior $\boldsymbol{\mu}$ encodes the distribution over
units weighted by scholarly output, and $\alpha$ controls the degree to
which influence propagates through the citation network rather than
being absorbed by that prior.
In the limit $\alpha\to 1$, equation~\eqref{eq:spectral} reduces to the
principal eigenvector problem
\begin{equation}
    \boldsymbol{\pi}^{\infty}
    \;=\;
    \mathbf{H}^\top\boldsymbol{\pi}^{\infty},
    \qquad
    \mathbf{1}^\top\boldsymbol{\pi}^{\infty}=1,
    \quad
    \boldsymbol{\pi}^{\infty}\ge 0,
    \label{eq:eig}
\end{equation}
whose solution is unique when $\mathbf{H}$ is primitive (irreducible and
aperiodic). \cite{gellerCitationInfluenceMethodology1978} shows that
$\boldsymbol{\pi}^{\infty}$ is then the stationary distribution of the
Markov chain with transition matrix $\mathbf{H}$, and that it coincides
with the \textit{influence weight} of \cite{pinskiCitationInfluenceJournal1976}
scaled by the reference count proportion $r_p / \sum_q r_q$.

The spectral ranking $\boldsymbol{\pi}$ reflects the total citation flow
received by each unit and therefore depends on unit size.
\emph{Prestige per work} \citep{palacios-huertaMeasurementIntellectualInfluence2004},
\begin{equation}
    \mathbf{v}(\alpha)
    \;=\;
    \mathbf{A}^{-1}\boldsymbol{\pi}(\alpha),
    \qquad
    \mathbf{A}=\mathrm{diag}(a_p),
    \label{eq:vstar}
\end{equation}
removes this dependence, where $a_p = |\{i:w_i\in s\}|$ for source $p=s$
and $a_p = \sum_i\omega_{iu}$ for institution $p=u$ via \eqref{eq:omega} are the work-counts,
$\boldsymbol{\pi}(\alpha)$ denotes the solution to \eqref{eq:spectral},
and $\boldsymbol{\pi}^{\infty}$ the limiting case \eqref{eq:eig}.
\cite{palacios-huertaMeasurementIntellectualInfluence2004} show that
$\mathbf{v}$ is uniquely characterised by several reasonable axiomatic
constraints \citep[but see][]{serranoMeasurementIntellectualInfluence2004}.
Note that $\mathbf{A}^{-1}$ must be applied after convergence; applying
it inside the iteration loop would iterate on $\mathbf{H}^\top\mathbf{A}$,
which is not norm-preserving and yields a fixed point with no clean
interpretation.

To compute $\boldsymbol{\pi}$ we define the modified matrix
\begin{equation}
	\mathbf{\tilde{H}}
	\;=\;
	\alpha\mathbf{H}+(1-\alpha)\boldsymbol{\mu}\mathbf{1}^\top,
	\label{eq:Hstar}
\end{equation}
so that the fixed-point condition
$\boldsymbol{\pi}=\mathbf{\tilde{H}}^{\top}\boldsymbol{\pi}$
is equivalent to equation~\eqref{eq:spectral}.
Because $\boldsymbol{\mu}>0$ and $0<\alpha<1$, every entry of $\mathbf{\tilde{H}}$
is strictly positive, ensuring it is a primitive row-stochastic matrix.
By the Perron--Frobenius theorem, $\mathbf{\tilde{H}}^{\top}$ therefore
has a unique stationary distribution $\boldsymbol{\pi}$, recovered by
power iteration
\begin{equation}
	\boldsymbol{\pi}^{(k+1)}
	\;=\;
	\mathbf{\tilde{H}}^{\top}\boldsymbol{\pi}^{(k)},
	\label{eq:power}
\end{equation}
converging to $\boldsymbol{\pi}$ from any strictly positive starting
vector $\boldsymbol{\pi}^{(0)}$.

\paragraph{Bipartite structure.}
With $\mathbf{m}=(0,1,1,0)$, equation~\eqref{eq:C} gives a block off-diagonal $\mathbf{H}$,
forming a bipartite directed graph between sources and institutions.
Since the product of row-stochastic matrices is row-stochastic, $\mathbf{H}^2$ is
block-diagonal and row-stochastic, with one-mode projection blocks
$\mathbf{M}_S=\mathbf{H}_{SI}\mathbf{H}_{IS}$ and $\mathbf{M}_I=\mathbf{H}_{IS}\mathbf{H}_{SI}$.
The bipartite structure makes $\mathbf{H}$ periodic with period~2, so the undamped
equation~\eqref{eq:eig} need not have a unique solution; $\tilde{\mathbf{H}}$ in
equation~\eqref{eq:Hstar} restores primitivity for any $\alpha\in(0,1)$.
Partitioning equation~\eqref{eq:spectral} and eliminating $\boldsymbol{\pi}_I$ yields
the Katz--Hubbell resolvent for sources,
\begin{equation}
  \bigl(\mathbf{I}-\alpha^2\mathbf{M}_S^\top\bigr)\boldsymbol{\pi}_S
  \;=\;
  (1-\alpha)\bigl(\boldsymbol{\mu}_S+\alpha\,\mathbf{H}_{IS}^\top\boldsymbol{\mu}_I\bigr),
  \label{eq:kh_source}
\end{equation}
where the effective damping is $\alpha^2$ and the effective prior transports institution
mass back to sources through $\mathbf{H}_{IS}^\top$. The institution ranking is then
recovered by a single matrix--vector product,
\begin{equation}
  \boldsymbol{\pi}_I
  \;=\;
  \alpha\,\mathbf{H}_{SI}^\top\boldsymbol{\pi}_S+(1-\alpha)\boldsymbol{\mu}_I.
  \label{eq:kh_inst}
\end{equation}

$\mathbf{C}(\chi,\mathbf{m},\rho)$ is a sparse matrix built from a
weighted edge list $\mathcal{E}$ accumulated over $\mathfrak{R}$ with
per-reference weight $\rho_i$, stored in CSR format, and
row-normalised to $\mathbf{H}$. The pipeline
$\mathfrak{R}\!\to\!\mathcal{E}\!\to\!\mathbf{C}\!\to\!\mathbf{H}
\!\to\!\boldsymbol{\pi}\!\to\!\mathbf{v}$
is implemented using \texttt{scipy.sparse.csr\_matrix} with parameter
entry points detailed in the Supplement.
